% !TeX root = ../main.tex
\section{Chân dung Khách hàng điển hình}

\subsection{Chân dung Người thuê thiết bị}

Hai persona này tập trung vào nhóm \textbf{Game thủ học sinh, sinh viên (Casual \& Student Gamers)} tại các khu vực đô thị và làng đại học lớn, đặc biệt là những cụm có mật độ sinh viên cao như Khu đô thị ĐHQG-HCM. Đây là tệp người thuê cốt lõi của Mutux trong giai đoạn đầu vì có nhu cầu trải nghiệm gaming gear cao, ngân sách mua mới còn hạn chế, thường sinh hoạt trong các cộng đồng trường/lớp/ký túc xá và có khả năng lan tỏa tự nhiên qua bạn bè, câu lạc bộ, phòng máy hoặc nhóm Discord. Việc tập trung vào nhóm này giúp nền tảng dễ đạt lượng người dùng hoạt động tối thiểu (critical mass) hơn so với việc tiếp cận nhiều phân khúc rời rạc ngay từ đầu.

\subsubsection{Persona 1: Minh Quân -- Sinh viên muốn trải nghiệm gaming gear tốt, mới ra mắt}

\begin{itemize}
    \item \textbf{Nhân khẩu học:} Nam, 20 tuổi, sinh viên năm hai ngành Công nghệ thông tin tại Trường Đại học Công nghệ Thông tin, ĐHQG-HCM. Sống tại ký túc xá ĐHQG, ngân sách cá nhân khoảng 2--4 triệu đồng/tháng, thường chi tiêu cho ăn uống, đi lại, game và một phần nhỏ cho thiết bị học tập/giải trí.
    \item \textbf{Bối cảnh sử dụng:} Quân chơi Valorant, CS2 và các game co-op cùng nhóm bạn 4--5 buổi vào tối hàng tuần tại ký túc xá. Quân muốn dùng thử chuột nhẹ, bàn phím cơ, tai nghe tốt hơn thiết bị phổ thông đang có và các mẫu gaming gear mới ra mắt để chơi thoải mái hơn, bắt kịp xu hướng trong nhóm bạn và biết sản phẩm nào hợp với mình trước khi mua.
    \item \textbf{Nỗi đau:} Không đủ ngân sách mua ngay một bộ gear 3--6 triệu đồng; các sản phẩm mới ra mắt thường có giá cao, vượt quá khả năng chi trả của sinh viên; sợ mua nhầm chuột không hợp tay hoặc switch bàn phím quá ồn trong phòng trọ; chỉ muốn trải nghiệm trong vài ngày/cuối tuần nhưng thị trường chủ yếu bán mới; lo ngại thuê/mượn qua người lạ vì không rõ tình trạng thiết bị và trách nhiệm khi hỏng hóc.
    \item \textbf{Động lực:} Muốn nâng trải nghiệm chơi game và setup cá nhân mà không phải đầu tư lớn; muốn thử thiết bị thực tế vài ngày trước khi quyết định mua; thích được trải nghiệm sớm sản phẩm mới và đổi mới gear theo xu hướng trong cộng đồng bạn bè. Vì ngân sách hạn chế, Quân đặc biệt quan tâm đến mô hình cọc linh hoạt, chẳng hạn cọc trả sau hoặc giảm áp lực cọc ban đầu cho người thuê có lịch sử tốt.
    \item \textbf{Hành vi tiếp cận:} Dễ gặp nhất khi chơi game chung, qua vòng tròn bạn bè, nhóm Discord nhỏ, Facebook group trường, ký túc xá và lời giới thiệu từ người quen đã từng thuê. Đây là persona phù hợp nhất cho MVP vì có mật độ tập trung cao, phản hồi nhanh và khả năng tạo word-of-mouth trong cộng đồng sinh viên.
\end{itemize}

\subsubsection{Persona 2: Nhật Khang -- Game thủ có thi đấu giải}

\begin{itemize}
    \item \textbf{Nhân khẩu học:} Nam, 22 tuổi, sinh viên năm cuối ngành Công nghệ thông tin tại Trường Đại học Khoa học Tự nhiên, ĐHQG-HCM, thường tham gia các giải Valorant/CS2 sinh viên và giải cộng đồng bán chuyên. Có đội chơi cố định, tự góp ngân sách cá nhân hoặc quỹ đội để đăng ký giải, thuê phòng máy và chuẩn bị thiết bị.
    \item \textbf{Bối cảnh sử dụng:} Nhu cầu của Khang tăng mạnh theo mùa giải. Trước các giải kéo dài 2--5 ngày, Khang cần chuột nhẹ, bàn phím phản hồi nhanh và tai nghe định vị âm thanh tốt để luyện tập và thi đấu ổn định. Ngoài mùa giải, nhu cầu thuê thấp hơn; trong mùa giải, thời gian chuẩn bị ngắn khiến việc có gear đúng mẫu, đúng thời điểm trở nên quan trọng.
    \item \textbf{Nỗi đau:} Mua gear chuẩn esports quá tốn kém nếu chỉ dùng nhiều trong một mùa giải hoặc một giải ngắn ngày; khó tiếp cận thiết bị đúng chuẩn thi đấu với chi phí hợp lý; mượn thiết bị qua người quen không ổn định và thiếu cam kết về thời gian giao nhận; cần giao nhận nhanh, đúng mẫu thiết bị và tình trạng kỹ thuật rõ ràng.
    \item \textbf{Động lực:} Muốn thuê ngắn ngày hoặc theo mùa giải thiết bị chuẩn esports, độ trễ thấp và tình trạng tốt; sẵn sàng trả phí cao hơn thuê thông thường nếu thiết bị giúp chuẩn bị thi đấu ổn định; quan tâm đến quy trình đặt cọc, bàn giao và hoàn trả minh bạch; có động lực giới thiệu dịch vụ cho đồng đội hoặc đội khác nếu Mutux đáp ứng tốt giai đoạn cao điểm mùa giải.
    \item \textbf{Hành vi tiếp cận:} Dễ gặp qua lịch giải đấu sinh viên, phòng máy, câu lạc bộ esports, admin cộng đồng game và đối tác tổ chức sự kiện. Persona này có mức độ nỗi đau rất cao theo mùa giải, phù hợp để \projectname{} thử nghiệm gói thuê ngắn ngày/theo mùa cho game thủ sinh viên có thi đấu giải và tạo trường hợp tham chiếu trong cộng đồng esports.
\end{itemize}

Nhìn chung, hai persona người thuê đều thuộc cùng một cụm khách hàng mục tiêu ban đầu: học sinh, sinh viên chơi game tại các khu vực tập trung đông người trẻ. Họ khác nhau về cường độ sử dụng và bối cảnh thuê, nhưng cùng chia sẻ nhu cầu trải nghiệm gear tốt với chi phí thấp hơn mua mới, khả năng phản hồi nhanh và khả năng giới thiệu dịch vụ cho bạn bè. Vì vậy, thay vì mở rộng ngay sang nhiều nhóm người thuê rời rạc, Mutux nên ưu tiên đạt độ phủ trong các cộng đồng trường đại học, ký túc xá, phòng máy và câu lạc bộ esports để tạo critical mass cho MVP.

% \begin{table}[H]
%     \centering
%     \caption{So sánh hai chân dung người thuê điển hình}
%     \label{tab:renter-personas-pa3}
%     \small
%     \begin{tabular}{p{0.22\textwidth}p{0.34\textwidth}p{0.34\textwidth}}
%         \toprule
%         \textbf{Tiêu chí} & \textbf{Minh Quân -- Sinh viên} & \textbf{Nhật Khang -- Game thủ thi đấu giải} \\
%         \midrule
%         Bối cảnh sử dụng & Thuê gear ngắn ngày để chơi cùng nhóm, trải nghiệm setup tốt hơn và thử trước khi mua. & Thuê gear chuẩn thi đấu trước và trong giải. \\
%         Mức độ nỗi đau & Cao do ngân sách mua mới hạn chế và sợ chọn sai thiết bị. & Rất cao trước giải; thiết bị lỗi có thể ảnh hưởng trực tiếp đến thành tích đội. \\
%         Khả năng chi trả & Trả được đơn nhỏ theo ngày; có thể chịu ảnh hưởng từ bạn bè/gia đình khi mua mới. & Tự trả hoặc góp cùng đội/CLB; ưu tiên gói thuê ngắn ngày. \\
%         Kênh tiếp cận & CLB esports, nhóm trường, ký túc xá, Discord, phòng máy. & Giải đấu sinh viên, phòng máy, CLB esports, admin cộng đồng game. \\
%         Vai trò với MVP & Tạo giao dịch đầu, phản hồi nhanh và lan truyền trong cộng đồng sinh viên. & Kiểm chứng nhu cầu thuê mùa giải và tạo trường hợp tham chiếu trong cộng đồng esports. \\
%         \bottomrule
%     \end{tabular}
% \end{table}

\subsection{Chân dung Người cho thuê Thiết bị}

Phía cung của Mutux không nên được hiểu là mọi cá nhân đang sở hữu một bộ gear đều sẵn sàng cho thuê. Khi các thiết bị phổ biến ngày càng dễ mua và giá thuê của một món đồ không còn cách quá xa giá mua, động lực cho thuê của cá nhân sẽ thấp; trong khi đó, chi phí kiểm tra, vệ sinh, giao nhận và rủi ro hư hỏng vẫn tồn tại. Vì vậy, ở giai đoạn đầu, Mutux ưu tiên các đối tác đã có sẵn thiết bị, không gian và tệp khách hàng. Với họ, cho thuê là một dịch vụ bổ sung để tăng trải nghiệm hoặc doanh thu, thay vì trở thành nguồn thu chính.

\subsubsection{Persona 3: Nam -- Chủ quán net/gaming center}

\begin{itemize}
    \item \textbf{Bối cảnh:} Nam vận hành một quán net hoặc gaming center quy mô nhỏ đến vừa quanh khu vực trường đại học, khu dân cư và các điểm tập trung game thủ. Quán đã có sẵn chuột, bàn phím, tai nghe, tay cầm hoặc một vài bộ máy dự phòng để phục vụ khách và các nhóm chơi theo giờ.
    \item \textbf{Nhu cầu và vấn đề:} Doanh thu chính vẫn đến từ giờ chơi, nhưng thiết bị dự phòng hoặc phụ kiện ít được sử dụng có thể bị khấu hao mà chưa tạo thêm giá trị. Nam cũng muốn giữ khách bằng trải nghiệm khác biệt, song không muốn tự xây dựng một quy trình cho thuê riêng, phải theo dõi đặt cọc, tình trạng thiết bị và giao nhận ngoài hoạt động của quán.
    \item \textbf{Động lực tham gia:} Có thêm doanh thu từ những thiết bị sẵn có; cho phép khách thuê nhanh phụ kiện khi quên đồ hoặc muốn thử một setup tốt hơn; hỗ trợ các đội game, câu lạc bộ và giải đấu nhỏ tổ chức tại quán. Mutux đóng vai trò mang thêm khách và xử lý phần đặt thuê, còn quán giữ quyền quyết định thiết bị nào được đưa lên hệ thống.
    \item \textbf{Rào cản:} Lo ngại thất lạc, hư hỏng, vệ sinh thiết bị và việc cho thuê làm ảnh hưởng đến số máy phục vụ khách tại giờ cao điểm. Do đó, Nam phù hợp với mô hình thử nghiệm phạm vi hẹp: vài phụ kiện có nhu cầu rõ ràng, nhận--trả tại quán và quy định bồi thường đơn giản.
    \item \textbf{Hành vi tiếp cận:} Có thể tiếp cận qua các quán net gần HCMUS, Thủ Đức, câu lạc bộ esports, admin cộng đồng game và đơn vị tổ chức giải sinh viên. Đây là persona cung cấp phù hợp nhất cho pilot vì đã có sẵn tài sản, nhân sự hỗ trợ và địa điểm nhận--trả.
\end{itemize}

\subsubsection{Persona 4: Bình -- Chủ shop gear quy mô nhỏ}

\begin{itemize}
    \item \textbf{Bối cảnh:} Bình điều hành một shop gear độc lập có showroom nhỏ hoặc kết hợp bán hàng online. Shop bán chuột, bàn phím, tai nghe, màn hình và phụ kiện gaming; một phần hàng được dùng để trưng bày, cho khách trải nghiệm, hoặc là hàng cũ/đổi trả còn sử dụng tốt. Các showroom gear vốn đã có vai trò cho khách xem và trải nghiệm sản phẩm trực tiếp, như mô hình showroom của GEARVN tại Thành phố Thủ Đức \cite{gearvnShowroom}.
    \item \textbf{Nhu cầu và vấn đề:} Biên lợi nhuận của các món phổ thông có thể không cao, còn khách hàng thường muốn thử cảm giác cầm chuột, độ nảy bàn phím hoặc chất âm trước khi mua. Tuy nhiên, Bình không muốn mọi sản phẩm trong kho đều phải chuyển sang cho thuê vì sẽ phát sinh kiểm kê, vệ sinh, hao mòn và rủi ro làm giảm giá trị bán lại.
    \item \textbf{Động lực tham gia:} Dùng Mutux như một kênh ``dùng thử trước khi mua'', biến hàng demo hoặc một số sản phẩm chọn lọc thành doanh thu phụ, đồng thời đưa khách thuê quay lại shop để mua phụ kiện phù hợp. Mục tiêu chính của Bình vẫn là bán hàng; khoản phí thuê và dữ liệu phản hồi chỉ là phần bổ trợ giúp tăng khả năng chốt đơn.
    \item \textbf{Rào cản:} Cần quy định rõ thiết bị nào được cho thuê, mức đặt cọc, trách nhiệm khi trầy xước và cách đối soát doanh thu. Shop chỉ nên đưa lên Mutux các thiết bị có giá trị trải nghiệm cao hoặc khó quyết định khi mua, thay vì các món giá rẻ mà khách có thể mua ngay.
    \item \textbf{Hành vi tiếp cận:} Phù hợp để tiếp cận trực tiếp qua các shop địa phương, cộng đồng gear, sự kiện gaming và đối tác tổ chức giải. Một gói thử nghiệm đơn giản (một vài sản phẩm, thời hạn ngắn, nhận--trả tại shop) sẽ giúp Bình đánh giá nhu cầu mà không phải thay đổi mô hình bán hàng hiện tại.
\end{itemize}

\subsubsection{Hàm ý cho mô hình cung ban đầu}

Hai persona trên cho thấy nguồn cung khả thi của Mutux nên bắt đầu từ đối tác kinh doanh có sẵn hoạt động vận hành, thay vì kỳ vọng vào người dùng cá nhân. Quán net tạo ra điểm nhận--trả và nhu cầu thuê tại chỗ; shop gear tạo ra điểm trải nghiệm và cơ hội chuyển đổi sang mua hàng. Các cá nhân sở hữu gear có thể được mở rộng ở giai đoạn sau, khi Mutux đã chứng minh quy trình đặt cọc, kiểm tra và xử lý hư hỏng đủ đơn giản để họ cảm thấy an toàn.
